% Causal-Epistemic Type Theory with Mechanized Metatheory
% Target: ICFP 2027
%
% The THEORY paper. Pure formalism + Lean proofs.
% No LOC counts, no benchmarks, no "self-hosted" marketing.
%
% Compile: pdflatex main && bibtex main && pdflatex main && pdflatex main

\documentclass[acmsmall,screen,review,anonymous]{acmart}

% Packages
\usepackage{mathpartir}
\usepackage{stmaryrd}
\usepackage{booktabs}
\usepackage{subcaption}
\usepackage{xspace}

% Macros
\newcommand{\sounio}{\textsc{Sounio}\xspace}
\newcommand{\know}[1]{\texttt{Knowledge}\langle #1 \rangle}
\newcommand{\knowfull}[4]{\texttt{Knowledge}\langle #1, #2, #3, #4 \rangle}
\newcommand{\cknow}[3]{\texttt{CausalKnowledge}\langle #1, #2, #3 \rangle}
\newcommand{\eps}{\epsilon}
\newcommand{\lean}{\textsc{Lean\,4}\xspace}

% Code listing style
\lstdefinelanguage{sounio}{
  keywords={fn, let, var, if, else, while, for, return, struct, enum, match,
            with, kernel, linear, type, use, effect, handle, perform, do},
  keywordstyle=\color{blue!60!black}\bfseries,
  ndkeywords={Knowledge, CausalKnowledge, CausalGraph, f64, i32, bool,
              IO, Mut, Alloc, GPU, Prob, Epistemic},
  ndkeywordstyle=\color{purple!60!black},
  sensitive=true,
  comment=[l]{//},
  commentstyle=\color{green!40!black}\itshape,
  stringstyle=\color{red!60!black},
  morestring=[b]",
}

\lstset{
  language=sounio,
  basicstyle=\ttfamily\small,
  frame=single,
  numbers=left,
  numberstyle=\tiny\color{gray},
  breaklines=true,
  captionpos=b,
  xleftmargin=2em,
  aboveskip=0.8em,
  belowskip=0.5em,
}

% ACM metadata
\setcopyright{acmcopyright}
\acmJournal{PACMPL}
\acmVolume{1}
\acmNumber{ICFP}
\acmArticle{1}
\acmYear{2027}
\acmMonth{9}
\acmDOI{10.1145/nnnnnnn.nnnnnnn}

\begin{document}

\title{Causal-Epistemic Type Theory: \\
Mechanized Metatheory for Uncertainty-Aware Causal Reasoning}

\author{Anonymous}
\affiliation{%
  \institution{Anonymous Institution}
  \country{Anonymous}
}
\email{anonymous@example.org}

\begin{abstract}
Scientific reasoning requires both \emph{epistemic judgment}---how
precisely is a quantity known?---and \emph{causal judgment}---what
would happen under intervention?
We present a type system that unifies these two concerns.
\emph{Epistemic types} $\know{T}$ attach measurement uncertainty and
provenance to values, with arithmetic that automatically propagates
uncertainty following the Guide to the Expression of Uncertainty in
Measurement (GUM, ISO/IEC~98-3).
\emph{Causal types} $\cknow{Y}{\eps}{\texttt{do}(X)}$ embed Pearl's
do-calculus into the type system, rejecting programs that compute
causal effects from non-identifiable queries.

We prove that the combined system is type-safe (progress and
preservation) and GUM-compliant (the computed uncertainty equals the
GUM combined standard uncertainty for all well-typed programs).
These results compose with linear types (four-modality lattice),
algebraic effects (row-polymorphic effect tracking), and units of
measure (dimensional analysis).

\textbf{All metatheory is mechanized in \lean: 731 definitions and
theorems with zero \texttt{sorry} and zero Mathlib dependency.}
To our knowledge, this is the first mechanized treatment of
Pearl's causal hierarchy in a type system context, and the first
type safety proof for a language combining epistemic uncertainty,
causal reasoning, linear resources, algebraic effects, and
dimensional analysis.
\end{abstract}

\begin{CCSXML}
<ccs2012>
<concept><concept_id>10011007.10011006.10011008.10011009.10011012</concept_id>
<concept_desc>Software and its engineering~Functional languages</concept_desc>
<concept_significance>500</concept_significance></concept>
<concept><concept_id>10003752.10010124.10010138</concept_id>
<concept_desc>Theory of computation~Program verification</concept_desc>
<concept_significance>500</concept_significance></concept>
<concept><concept_id>10003752.10010124.10010131</concept_id>
<concept_desc>Theory of computation~Type theory</concept_desc>
<concept_significance>500</concept_significance></concept>
</ccs2012>
\end{CCSXML}

\ccsdesc[500]{Software and its engineering~Functional languages}
\ccsdesc[500]{Theory of computation~Program verification}
\ccsdesc[500]{Theory of computation~Type theory}

\keywords{epistemic types, causal inference, do-calculus, type safety,
mechanized metatheory, GUM, Lean~4, uncertainty quantification}

\maketitle

% ============================================================================
\section{Introduction}
\label{sec:intro}
% ============================================================================

Every measured quantity carries uncertainty.
A laboratory balance reports mass as $m = 10.023 \pm 0.002\,\text{g}$;
a clinical assay yields serum creatinine as $\text{SCr} = 1.2 \pm 0.1\,\text{mg/dL}$;
an ODE solver accumulates numerical error proportional to step count.
When these quantities enter computation---drug dosing calculations,
climate projections, neural network training---their uncertainties must
propagate correctly through every arithmetic operation,
or the results are scientifically meaningless.

But uncertainty alone is insufficient.
A clinician does not merely ask ``what is this patient's creatinine?''
but ``what \emph{would} the creatinine be if we changed the dosing
regimen?''---a causal question that cannot be answered from
observational data alone without structural assumptions.
Current programming languages offer no type-level support for either
concern: uncertainty is tracked (if at all) by runtime libraries
that provide no static guarantees~\cite{lebigot2024uncertainties,
giordano2016measurements}, and causal reasoning is entirely
outside the type system.

We present a type system that makes both epistemic and causal
reasoning first-class type-level concerns, with the following
contributions:

\begin{enumerate}
  \item A type system for \emph{epistemic values} with formal typing
        rules, operational semantics, and proofs of progress,
        preservation, and GUM compliance soundness
        (\S\ref{sec:epistemic}).

  \item A type system for \emph{causal reasoning} that embeds Pearl's
        do-calculus~\cite{pearl2009causality} into the type system,
        rejecting programs that compute causal effects from
        non-identifiable queries (\S\ref{sec:causal}).

  \item The \emph{composition} of epistemic and causal types with
        linear types, algebraic effects, and units of measure,
        yielding a unified type system for scientific reasoning
        (\S\ref{sec:composition}).

  \item \textbf{Full mechanization in \lean}: 731 definitions and
        theorems covering type safety, epistemic confidence lattices,
        causal type theory with do-calculus, linear types with
        four-modality lattice, algebraic effect rows, and units of
        measure---all with zero \texttt{sorry} and zero Mathlib
        dependency (\S\ref{sec:mechanization}).
\end{enumerate}

% ============================================================================
\section{Epistemic Types}
\label{sec:epistemic}
% ============================================================================

% TODO: Extract and refine from popl2027/main.tex sections 2.1--2.3
% Focus on: Syntax, Typing Rules, Operational Semantics
% Include: T-Measure, T-Add, T-Mul, T-Transform, T-Extract
% Include: E-Add, E-Mul, E-Extract, E-Beta, E-LetP

\textbf{[Content from POPL S2: Epistemic Type System --- to be extracted]}

% ============================================================================
\section{Causal Types}
\label{sec:causal}
% ============================================================================

% TODO: Expand from popl2027/main.tex section 5.4 (Causal Types)
% This section becomes MUCH longer --- it's now a core contribution, not an extension.
%
% Subsections:
% 3.1 Structural Causal Models (syntax + typing)
% 3.2 Do-Operator at the Type Level (T-Do rule)
% 3.3 Identifiability as a Type Constraint
% 3.4 d-Separation and Conditional Independence
% 3.5 Back-Door and Front-Door Criteria
% 3.6 Causal-Epistemic Interaction

\textbf{[Content: Expand causal types from POPL S5.4 into full section]}

% ============================================================================
\section{Metatheory}
\label{sec:metatheory}
% ============================================================================

% TODO: Extract from popl2027/main.tex section 3
% Theorems: Determinism, Progress, Preservation, Type Safety, GUM Soundness
% ADD: Causal-specific theorems (identifiability soundness, d-sep monotonicity)

\textbf{[Content from POPL S3: Metatheory --- to be extracted]}

% ============================================================================
\section{Composition with Other Type Disciplines}
\label{sec:composition}
% ============================================================================

% TODO: Extract from popl2027/main.tex section 5 (Extensions)
% Linear types, algebraic effects, units of measure
% Key point: all compose cleanly with epistemic + causal

\textbf{[Content from POPL S5: Extensions --- to be extracted]}

% ============================================================================
\section{Mechanization in Lean 4}
\label{sec:mechanization}
% ============================================================================

% TODO: Extract from popl2027/main.tex section 4
% Table of modules, proof architecture, novel results
% EXPAND: Causal theorem inventory (106 theorems)

\textbf{[Content from POPL S4: Mechanization --- to be extracted]}

% ============================================================================
\section{Motivating Example}
\label{sec:example}
% ============================================================================

% A SHORT (10-15 line) example showing epistemic + causal types in action.
% E.g., pharmacokinetic intervention: "what would the drug concentration
% be if we changed the dosing regimen?" with uncertainty bounds.
% This is NOT a benchmark or evaluation --- just a taste.

\textbf{[One concise motivating example]}

% ============================================================================
\section{Related Work}
\label{sec:related}
% ============================================================================

% TODO: Extract and expand from popl2027/main.tex section 7
% ADD: More causal inference references (Bareinboim, Tian, Shpitser)
% ADD: Comparison to probabilistic programming for causal inference

\textbf{[Content from POPL S7: Related Work --- to be expanded]}

% ============================================================================
\section{Conclusion}
\label{sec:conclusion}
% ============================================================================

\textbf{[To be written]}

\bibliographystyle{ACM-Reference-Format}
\bibliography{references}

\end{document}
